% Referee report 3, Major 1 (finiteness: contact kernels, contact cancellation, transition region of (C10))
% and Major 6 (conditional convergence: (C3), (C7), (F.12), the two integrals of (6.15)).
% Scripts: paper/notes/scripts/ref3_finite/ (python -I; numpy/scipy/mpmath).  Wrapper: ref3_finiteness_wrapper.tex.
\def\v#1{\mathbf{#1}}
\providecommand{\Kv}{\mathcal K}
\providecommand{\mf}[1]{\mathbb{#1}}
\section{Finiteness of the contact terms and conditional convergence (referee 3, Major 1 and Major 6)}\label{sec:ref3fin}

\noindent\textbf{Result in one paragraph.}
(a) Every contact kernel of the $(2\rho)$ equations (31 printed lines of Eqs.~(sIIII), (sIV1), (sIII1), (sIIIII), (sIIIIII1), (sIIIV), (sIIIVI), decomposed into 47 kernel$\times$colour units, plus the two lines of (sIIb)) was power-counted at \emph{every} coincidence of its integration points (1222 clusters of two to six points, 705 of them containing the merged charge point) and, for a local two-point function $n$, at the triple coincidence of the charges (329 clusters). With the colour factors as printed, every line is \emph{absolutely} integrable at every coincidence; no angular average, principal value at a point, or cancellation between lines is needed locally, except the subtraction $\mathfrak K^{(y')}_\delta\mathcal L(y')-\mathfrak K^{(z)}_\delta\mathcal L(z)$ at $\v y'\to\v z$ which the paper already prescribes. The referee's kernel $K(x,x,z)=1/(\v x-\v z)^2$ occurs in exactly one place, $\mathfrak K^{(z)}_\delta|_{x'\to x}$ of Rows III.III, III.V and III.VI; there the kernel alone is not integrable at $\v z\to\v x$ (and at $\v x,\v y',\v z\to$ one point), but the colour factor of each of the three lines vanishes at $\v z=\v x$, identically in the Wilson lines (by $\mf F(a,a,c)=N_c\mf D(a,c)$), so the line is integrable for every target whose Wilson lines are H\"older continuous; in all other lines the kernel is integrable by itself. The large-distance tails of the contact lines cancel pair by pair as at separate points (operator identity; checked numerically for pairs A and D in all contacts). (b) The contact terms do \emph{not} cancel among the rows at the integrand level: lines with the same kernel exist only within Group~III (Rows III.III, III.V, III.VI) and between Rows II.I and III.II, and the sums of their colour factors are non-zero (in II.I$+$III.II the quadrupole $\mf Q$ terms cancel, dipole and $\mf F$ terms remain); a cancellation against the three-$\rho$ evolution contact terms is excluded because those multiply $\ln(\vee/k^+)$ or $\ln(k^+/\Lambda)$. Finiteness does not need a cancellation. (c) The transition region of (C10) contributes nothing at the order of the tail: after the measured-gluon Fourier transform the large-$\v z$ expansion of every two-gluon amplitude is fixed by the inner regions only (projectile, pair, and collinear configuration), the region $|\v w|\sim|\v w-\v z|\sim|\v z|$ being smooth; verified numerically on the full $\v w$ plane (residual $\propto\ln Z/Z$ resp.\ $1/Z$ for $t=1$, $0.3$, $Z\le80$). (d) None of (C3), (6.15), (F.12), (C7) is a prescription that could shift the result: (C3), (6.15) and (F.12) are Fourier integrals at the measured momentum, which are unique (independent of the regulator, the order radial/angular and the cut-off shape) by a standard lemma, after which every remaining integral is absolutely convergent; for (C7) the only non-absolutely convergent piece is $\v k$-independent and is annihilated by $\int d\kappa\,e^{i\kappa}[\pi\delta(\kappa)+i\mathcal P\frac1\kappa]=0$. Each statement was checked numerically with several regulators (agreement to $\le10^{-3}$ at $\epsilon=10^{-4}$, error $\propto\epsilon$; (F.12): two orders agree to $5\times10^{-11}$). One genuinely conditionally convergent \emph{local} point remains, not on the referee's list: the collinear point $\v z\to\v x$ of $\mathfrak K^{\rm V}_{2b}$ (Row I.V), a pure second harmonic $a_2\cos2(\varphi-\varphi_0)/r^2$; its value is fixed by the rotation invariance of the regularisation of the accompanying $\frac{11}{6}$ subtraction and is tied to the $d_\perp$ question of Major~4 (Sec.~\ref{sec:r3f-IV}). \textbf{No correction of the finite terms is required.} Drop-in text is in Sec.~\ref{sec:r3f-dropin}.

\subsection{Setup and method}\label{sec:r3f-setup}
\noindent A $(2\rho)$ line is $\mathcal C_{\rm f}(N_c^2-1)^{-1}\int(\dots)\,\mathfrak K|_{q\to p}\;\mathcal S\;n(p,r)$: a $c$-number kernel with the charge point $q$ merged into $p$, a singlet colour factor $\mathcal S$ (combination of $\mf D,\mf F,\mf Q$ at the integration points) and the two-point function $n$ of the remaining charges $p,r$. A coincidence is a cluster $S$ of integration points approaching one point at scale $\lambda$ with all other points fixed; with $|\mathfrak K\,\mathcal S|\sim\lambda^{\omega(S)}$ (up to logarithms), the line is absolutely integrable at $S$ iff $\omega(S)>-d_S$, $d_S=2(|S|-1)$ (one less if $S\ni\v z,\v y'$ in a $\delta(s)$ term, where $\v y'-\v z\perp\v k$). For kernels that are products of homogeneous functions of point differences, logarithms of ratios and bounded functions of $s$, absolute integrability at all clusters is sufficient for local integrability (Hepp sectors / Weinberg power counting). Principal values $\mathcal P\frac1s$ act on a line, not a point; they are treated in Sec.~\ref{sec:r3f-C7}.

\noindent\emph{Target.} The colour factors were evaluated with smooth adjoint $SU(3)$ fields $U(\v x)=\exp[A^a(\v x)f^a]$, $A^a$ a sum of plane waves (two different fields: 4 modes, amplitude $1.2$; 6 modes, amplitude $2$, twice the gradients), non-uniform at large distance. Where the colour factor is needed (below) it enters only through its vanishing at a coincidence; for a target with H\"older continuous Wilson lines, $|U(\v z)-U(\v x)|\le c|\v z-\v x|^\gamma$ with any $\gamma>0$ (e.g.\ the McLerran--Venugopalan target, $\gamma<1$), the margins below stay positive.

\noindent\emph{Numerics} [\texttt{lib.py}, \texttt{lines.py}, \texttt{scan.py}, \texttt{run\_variants.py}]. Kernels transcribed from Eqs.~(eq:Mij), (sGIkernels), (sGIIIkernels), (app\_III1a); colour factors parsed from the printed $(2\rho)$ equations (the invariants satisfy $\mf D(a,a)=N_c^2-1$, $\mf F(a,a,c)=N_c\mf D(a,c)$, $\mf Q(a,a;c,d)=N_c\mf D(c,d)$, $\mf Q(a,b;b,d)=\mf F(a,b,d)$ to $10^{-14}$). For each line and cluster: 3 generic base configurations, 16 random orientations, $\lambda=10^{-2}\dots10^{-5}$, $\omega$ from the last two values (and the previous pair as a check); colour included and colour set to one. The $\delta(s)$ terms were scanned with $\v y'$ on the line $\v z+u\,\hat{\v k}_\perp$ ($\delta$ form) and, separately, with the factor $1/s$ and $\v y'$ free ($\mathcal P$ form). Two clusters of $\mathfrak K^{\rm III}_z$ at $\v z\to\v w$ lose all digits in double precision; they were redone with 60 digits [\texttt{mp\_Mij.py}]: $\lambda\max|\mathfrak K|/\ln(1/\lambda)=0.444,0.419,0.412,0.409$ for $\lambda=10^{-3},10^{-5},10^{-7},10^{-9}$, i.e.\ $\ln(1/r)/r$, integrable.

\subsection{(a) Power counting of the contact kernels}\label{sec:r3f-a}
\noindent\textbf{Results table.} One row per line (kernel unit, contact). Column 4: most singular cluster with the colour factor ($S$, $d_S$, $\omega$); column 5: most singular cluster of the kernel alone ($S$, $d_S$, $\omega_0$); ``integrable'': all 1222 clusters of the line (also those not containing the merged point) with $\omega>-d_S$; ``reason'': ``kernel'' if the kernel alone is integrable at every cluster, otherwise the clusters where the vanishing colour factor is needed; ``local $n$'': the same with the third charge merged as well. The minimal margin $\omega+d_S$ over all lines and clusters is $0.90$ (a $\ln(1/r)/r$ at $d_S=2$); with the second field the exponents agree to $0.05$ (825 clusters compared).
{\footnotesize
\begin{center}
\renewcommand{\arraystretch}{1.12}\setlength{\tabcolsep}{2pt}
\begin{tabular}{|l|l|c|l|l|c|l|c|}\hline
row & kernel & contact & with colour: $S$, $d_S$, $\omega$ & kernel alone: $S$, $d_S$, $\omega_0$ & integrable & reason & local $n$\\\hline
I.III & $\mathfrak K^{\rm III}_z$ & $x'\to x$ & $\{x,w\}$, $2$, $-1.1$ & $\{x,w\}, 2$, $-1.1$ & yes & kernel & yes\\
I.III & $\mathfrak K^{\rm III}_z$ & $y\to x$ & $\{x,w\}$, $2$, $-1.1$ & $\{x,z\}, 2$, $-1.1$ & yes & kernel & yes\\
I.III & $\mathfrak K^{\rm III}_z$ & $x'\to y$ & $\{y,w'\}$, $2$, $-0.0$ & $\{y,z\}, 2$, $-1.0$ & yes & kernel & yes\\
I.III & $\mathfrak K^{\rm III}_w$ & $x'\to x$ & $\{x,z\}$, $2$, $-1.0$ & $\{x,z\}, 2$, $-1.0$ & yes & kernel & yes\\
I.III & $\mathfrak K^{\rm III}_w$ & $y\to x$ & $\{x,z\}$, $2$, $-1.0$ & $\{x,z\}, 2$, $-1.0$ & yes & kernel & yes\\
I.III & $\mathfrak K^{\rm III}_w$ & $x'\to y$ & $\{y,w'\}$, $2$, $-0.0$ & $\{y,w\}, 2$, $-1.0$ & yes & kernel & yes\\
I.V & $\mathfrak K^{\rm V}_1$ & $x'\to y$ & $\{y,z\}$, $2$, $-1.0$ & $\{y,w'\}, 2$, $-1.0$ & yes & kernel & yes\\
I.V & $\mathfrak K^{\rm V}_1$ & $x'\to y'$ & $\{y',x\}$, $2$, $-1.0$ & $\{y',w'\}, 2$, $-1.0$ & yes & kernel & yes\\
II.I & $\mathfrak K_{\mathcal T}(w,w')$ & $y\to x$ & $\{x,z\}$, $2$, $-1.1$ & $\{x,z\}, 2$, $-1.1$ & yes & kernel & yes\\
II.I & $\mathfrak K_{\mathcal T}(w,w')$ & $x'\to x$ & $\{x,w'\}$, $2$, $-1.0$ & $\{x,w'\}, 2$, $-1.0$ & yes & kernel & yes\\
II.I & $\mathfrak K_{\mathcal T}(w',w)$ & $x'\to x$ & $\{x,w\}$, $2$, $-1.0$ & $\{x,w\}, 2$, $-1.0$ & yes & kernel & yes\\
II.I & $\mathfrak K_{\mathcal T}(w',w)$ & $y\to x$ & $\{x,z\}$, $2$, $-1.1$ & $\{x,z\}, 2$, $-1.1$ & yes & kernel & yes\\
III.II & $\mathfrak K_{\mathcal T}(w,w')$ & $x'\to y$ & $\{y,w'\}$, $2$, $-1.0$ & $\{y,w'\}, 2$, $-1.0$ & yes & kernel & yes\\
III.II & $\mathfrak K_{\mathcal T}(w,w')$ & $x'\to x$ & $\{x,w'\}$, $2$, $-1.0$ & $\{x,w'\}, 2$, $-1.0$ & yes & kernel & yes\\
III.II & $\mathfrak K_{\mathcal T}(w,w')$ & $y\to x$ & $\{x,w\}$, $2$, $-0.1$ & $\{x,z\}, 2$, $-1.1$ & yes & kernel & yes\\
III.III & $\mathfrak K_{\widehat{\mathcal E}}$ & $x'\to y$ & $\{y,w\}$, $2$, $-1.0$ & $\{y,w\}, 2$, $-1.0$ & yes & kernel & yes\\
III.III & $\mathfrak K_\Psi$ & $x'\to y$ & $\{y,w\}$, $2$, $-1.0$ & $\{y,y'\}, 2$, $-1.0$ & yes & kernel & yes\\
III.III & $\mathfrak K_\delta^{(y')}{+}\mathfrak K_\delta^{(z)}$, $\delta$ & $x'\to y$ & $\{y,w\}$, $2$, $-1.0$ & $\{y,z\}, 2$, $-1.0$ & yes & kernel & yes\\
III.III & $\mathfrak K_\delta^{(y')}{+}\mathfrak K_\delta^{(z)}$, $\mathcal P$ & $x'\to y$ & $\{y,w\}$, $2$, $-1.0$ & $\{y,y',z\}, 4$, $-3.0$ & yes & kernel & yes\\
III.III & $\mathfrak K_{\widehat{\mathcal E}}$ & $x'\to x$ & $\{x,z\}$, $2$, $-1.0$ & $\{x,y',z\}, 4$, $-3.1$ & yes & kernel & yes\\
III.III & $\mathfrak K_\Psi$ & $x'\to x$ & $\{x,z\}$, $2$, $-1.0$ & $\{x,y'\}, 2$, $-1.0$ & yes & kernel & yes\\
III.III & $\mathfrak K_\delta^{(y')}{+}\mathfrak K_\delta^{(z)}$, $\delta$ & $x'\to x$ & $\{x,z\}$, $2$, $-1.0$ & $\{x,z\}, 2$, $-2.0$ & yes & colour $=0$ at \{x,z\}, \{x,y',z\} & yes\\
III.III & $\mathfrak K_\delta^{(y')}{+}\mathfrak K_\delta^{(z)}$, $\mathcal P$ & $x'\to x$ & $\{x,y',z\}$, $4$, $-3.0$ & $\{x,y',z\}, 4$, $-4.0$ & yes & colour $=0$ at \{x,z\}, \{x,y',z\} & yes\\
III.V & $\mathfrak K_{\widehat{\mathcal E}}$ & $x'\to x$ & $\{x,z\}$, $2$, $-1.0$ & $\{x,y',z\}, 4$, $-3.1$ & yes & kernel & yes\\
III.V & $\mathfrak K_\Psi$ & $x'\to x$ & $\{x,z\}$, $2$, $-1.0$ & $\{x,z\}, 2$, $-1.0$ & yes & kernel & yes\\
III.V & $\mathfrak K_\delta^{(y')}{+}\mathfrak K_\delta^{(z)}$, $\delta$ & $x'\to x$ & $\{x,z\}$, $2$, $-1.0$ & $\{x,z\}, 2$, $-2.0$ & yes & colour $=0$ at \{x,z\}, \{x,y',z\} & yes\\
III.V & $\mathfrak K_\delta^{(y')}{+}\mathfrak K_\delta^{(z)}$, $\mathcal P$ & $x'\to x$ & $\{x,y',z\}$, $4$, $-3.0$ & $\{x,z\}, 2$, $-2.0$ & yes & colour $=0$ at \{x,z\}, \{x,y',z\} & yes\\
III.V & $\mathfrak K_{\widehat{\mathcal E}}$ & $y\to x$ & $\{x,z\}$, $2$, $-1.0$ & $\{x,z\}, 2$, $-1.0$ & yes & kernel & yes\\
III.V & $\mathfrak K_\Psi$ & $y\to x$ & $\{x,z\}$, $2$, $-1.0$ & $\{x,w\}, 2$, $-1.0$ & yes & kernel & yes\\
III.V & $\mathfrak K_\delta^{(y')}{+}\mathfrak K_\delta^{(z)}$, $\delta$ & $y\to x$ & $\{x,z\}$, $2$, $-1.0$ & $\{x,w\}, 2$, $-1.0$ & yes & kernel & yes\\
III.V & $\mathfrak K_\delta^{(y')}{+}\mathfrak K_\delta^{(z)}$, $\mathcal P$ & $y\to x$ & $\{x,z\}$, $2$, $-1.0$ & $\{x,z\}, 2$, $-1.0$ & yes & kernel & yes\\
III.V & $\mathfrak K_{\widehat{\mathcal E}}$ & $x'\to y$ & $\{y,y'\}$, $2$, $-0.0$ & $\{y,w\}, 2$, $-1.0$ & yes & kernel & yes\\
III.V & $\mathfrak K_\Psi$ & $x'\to y$ & $\{y,y'\}$, $2$, $-0.0$ & $\{y,y'\}, 2$, $-1.0$ & yes & kernel & yes\\
III.V & $\mathfrak K_\delta^{(y')}{+}\mathfrak K_\delta^{(z)}$, $\delta$ & $x'\to y$ & $\{y,w\}$, $2$, $-0.0$ & $\{y,w\}, 2$, $-1.0$ & yes & kernel & yes\\
III.V & $\mathfrak K_\delta^{(y')}{+}\mathfrak K_\delta^{(z)}$, $\mathcal P$ & $x'\to y$ & $\{y,z\}$, $2$, $-0.0$ & $\{y,y',z\}, 4$, $-3.0$ & yes & kernel & yes\\
III.VI & $\mathfrak K_{\widehat{\mathcal E}}$ & $x'\to y$ & $\{y,w\}$, $2$, $-1.0$ & $\{y,y'\}, 2$, $-1.0$ & yes & kernel & yes\\
III.VI & $\mathfrak K_\Psi$ & $x'\to y$ & $\{y,w\}$, $2$, $-1.0$ & $\{y,w\}, 2$, $-1.0$ & yes & kernel & yes\\
III.VI & $\mathfrak K_\delta^{(y')}{+}\mathfrak K_\delta^{(z)}$, $\delta$ & $x'\to y$ & $\{y,w\}$, $2$, $-1.0$ & $\{y,z\}, 2$, $-1.0$ & yes & kernel & yes\\
III.VI & $\mathfrak K_\delta^{(y')}{+}\mathfrak K_\delta^{(z)}$, $\mathcal P$ & $x'\to y$ & $\{y,w\}$, $2$, $-1.0$ & $\{y,y',z\}, 4$, $-3.0$ & yes & kernel & yes\\
III.VI & $\mathfrak K_{\widehat{\mathcal E}}$ & $y\to x$ & $\{x,w\}$, $2$, $-1.0$ & $\{x,w\}, 2$, $-1.0$ & yes & kernel & yes\\
III.VI & $\mathfrak K_\Psi$ & $y\to x$ & $\{x,w\}$, $2$, $-1.0$ & $\{x,w\}, 2$, $-1.0$ & yes & kernel & yes\\
III.VI & $\mathfrak K_\delta^{(y')}{+}\mathfrak K_\delta^{(z)}$, $\delta$ & $y\to x$ & $\{x,w\}$, $2$, $-1.0$ & $\{x,w\}, 2$, $-1.0$ & yes & kernel & yes\\
III.VI & $\mathfrak K_\delta^{(y')}{+}\mathfrak K_\delta^{(z)}$, $\mathcal P$ & $y\to x$ & $\{x,w\}$, $2$, $-1.0$ & $\{x,w\}, 2$, $-1.0$ & yes & kernel & yes\\
III.VI & $\mathfrak K_{\widehat{\mathcal E}}$ & $x'\to x$ & $\{x,y'\}$, $2$, $-0.0$ & $\{x,y',z\}, 4$, $-3.1$ & yes & kernel & yes\\
III.VI & $\mathfrak K_\Psi$ & $x'\to x$ & $\{x,y'\}$, $2$, $-0.0$ & $\{x,z\}, 2$, $-1.0$ & yes & kernel & yes\\
III.VI & $\mathfrak K_\delta^{(y')}{+}\mathfrak K_\delta^{(z)}$, $\delta$ & $x'\to x$ & $\{x,y',z\}$, $3$, $-1.0$ & $\{x,z\}, 2$, $-2.0$ & yes & colour $=0$ at \{x,z\}, \{x,y',z\} & yes\\
III.VI & $\mathfrak K_\delta^{(y')}{+}\mathfrak K_\delta^{(z)}$, $\mathcal P$ & $x'\to x$ & $\{x,y'\}$, $2$, $-0.0$ & $\{x,y',z\}, 4$, $-4.0$ & yes & colour $=0$ at \{x,z\}, \{x,y',z\} & yes\\
I.II $N_{2b}$ & $\mathfrak K^{\rm II}_{2b}$ & $x'\to x$, $x'\to y$ & \multicolumn{2}{l|}{analytic: $\{p,w'\}$: $-1$; $\{x,w,w'\}$: $-2$ (log); $\{x,y,w,w'\}$: $-2$} & yes & kernel & yes\\\hline
\end{tabular}
\end{center}}
\noindent($\delta$ and $\mathcal P$ denote the two parts of $\mathfrak K^{(y')}_\delta+\mathfrak K^{(z)}_\delta$, each line with its own colour factor, $\mathcal S(\v y')$ for $\mathfrak K^{(y')}_\delta$ and $\mathcal S(\v z)$ for $\mathfrak K^{(z)}_\delta$, as printed; in the $\mathcal P$ form a cluster with $\v z,\v y'$ carries the factor $1/s$. The entry ``$-0.0$'' means a bounded line at the most singular cluster.)

\noindent\textbf{Where the kernel alone is not integrable: the referee's $K(x,x,z)$.} Of the 1222 clusters, the kernel alone fails ($\omega_0=-d_S$) in exactly twelve: $\{\v x,\v z\}$ ($\omega_0=-2$) and $\{\v x,\v y',\v z\}$ ($\omega_0=-3$ in the $\delta$ form, $d=3$; $-4$ in the $\mathcal P$ form, $d=4$) of the line $\mathfrak K^{(z)}_\delta|_{x'\to x}$ in Rows III.III, III.V, III.VI (two forms each). With $x'=x$,
\begin{align}
\mathfrak K^{(z)}_\delta\big|_{x'\to x}=-\Big[\pi\delta(s)-i\mathcal P\frac1s\Big]\,\Kv(\v y'-\v z)\cdot\Kv(\v y-\v w)\;K(\v x,\v x,\v z)\,,\qquad K(\v x,\v x,\v z)=\frac{1}{(\v x-\v z)^2}\,.
\end{align}
Its colour factors (fourth line of (sIIIIII1), fourth of (sIIIV), sixth of (sIIIVI)) are
\begin{align}
\mathcal S_{\rm III.III}&=-N_c\mf D(x,y)-N_c\mf D(z,w)+\mf F(x,y,z)+\mf F(x,z,w)\,,\notag\\
\mathcal S_{\rm III.V}&=-N_c\mf D(x,y)+N_c\mf D(x,w)+\mf F(x,y,z)-\mf F(x,z,w)\,,\notag\\
\mathcal S_{\rm III.VI}&=N_c\mf D(x,y)+N_c\mf D(y,z)-2\mf F(x,y,z)\,,
\end{align}
and $\mf F$ is totally symmetric with $\mf F(a,a,c)=N_c\mf D(a,c)$, so at $\v z=\v x$: $\mathcal S_{\rm III.III}=-N_c\mf D(x,y)-N_c\mf D(x,w)+N_c\mf D(x,y)+N_c\mf D(x,w)=0$, $\mathcal S_{\rm III.V}=0$, $\mathcal S_{\rm III.VI}=N_c\mf D(x,y)+N_c\mf D(y,x)-2N_c\mf D(x,y)=0$, for every Wilson-line configuration. Hence $|\mathcal S|\le c\,|\v z-\v x|^\gamma$ for H\"older continuous lines, and $\int d^2x\,|\v x-\v z|^{\gamma-2}<\infty$. (Numerically $\omega=-1$ for III.III, III.V and $\omega=0$ for III.VI: there $\mathcal S$ vanishes quadratically for smooth lines, since $N_c\partial\mf D-2\partial\mf F=0$ at the coincidence by $f^{lqp}f^{lmn}f^{apn}=\frac{N_c}2f^{aqm}$.) At the triple coincidence $\{\v x,\v y',\v z\}$ the two lines of the subtraction enter together, $\mathcal D_s\,\Kv(\v y'-\v z)\cdot\Kv(\v y-\v w)\,\Kv(\v x-\v z)\cdot[\Kv(\v x-\v y')\,\mathcal S(\v y')-\Kv(\v x-\v z)\,\mathcal S(\v z)]$, $\mathcal D_s=\pi\delta(s)-i\mathcal P\frac1s$; the kernel is $\lambda^{-3}$ on the line $s=0$ ($d=3$), and both colour factors vanish at $\v x=\v y'=\v z$ (e.g.\ III.III: $\mathcal S(\v y')=-N_c\mf D(x,y)+\mf F(x,y,y')+\mf F(x,z,w)-\mf F(y',z,w)\to-N_c\mf D(z,y)+N_c\mf D(z,y)+N_c\mf D(z,w)-N_c\mf D(z,w)=0$), giving $\lambda^{-3+\gamma}$: integrable. The referee's example is therefore correct as a statement about the kernel and does not make the term divergent: the term is the commutator remnant of a colour structure that vanishes when the soft gluon sits on the charge.

\noindent\textbf{All other lines.} In the remaining 1210 clusters the kernel alone is integrable, with margin $\ge0.90$; the most singular behaviours are $\ln(1/r)/r$ ($\mathfrak K^{\rm III}_z$ at $\{\v x,\v w\}$, $\{\v z,\v w\}$; $\mathfrak K_{\mathcal T}|_{y\to x}$ at $\{\v x,\v z\}$) and $1/r$, and at triple clusters $\lambda^{-3}$ or weaker in $d=4$. The cancellation that makes $\mathfrak K^{\rm III}_z|_{y\to x}$ and $\mathfrak K_{\mathcal T}|_{y\to x}$ integrable at $\v z\to\v x$ (where $\Kv^j(\v x-\v z)$ is contracted with terms $\propto(\v x-\v z)^j/(\v x-\v z)^2$) is local in the kernel: in $\mathcal T^{ij}\Kv^j(\v x-\v z)$ the two $1/(\v x-\v z)^2$ terms combine into $[\Kv^i(\v z-\v w)-\Kv^i(\v x-\v w)]/(\v x-\v z)^2=O(1/|\v x-\v z|)$, and in $\mathcal M^{ij}\Kv^j(\v x-\v z)$ the terms $\ln\frac{B}{A}B^{-1}[-a^i/A+c^iA/(C(A-B))]$ combine into $O(\ln B/|\v x-\v z|)$ ($a=x-w$, $b=x-z$, $c=z-w$). No colour factor is needed for these.

\noindent\emph{Row I.II, $N_{2b}$} (not scanned; the coordinate form of $\mathfrak K^{\rm II}_{2b}$ is not available). The contact only sets $P^i=\Kv^i(\v x-\v w')$ or $\Kv^i(\v y-\v w')$; $\mathfrak K^{\rm II}_{2b}(\v x,\v y,\v w)$ does not depend on $\v w'$. $\mathcal R^i/\v q^2$ is exactly homogeneous of degree $-3$ in $(\v p,\v q)$ (its logarithms are of ratios), so $\mathfrak K^{\rm II}_{2b}$ is homogeneous of degree $-1$ under a joint scaling of $\v x-\v w,\v y-\v w$. Its sub-singularities: at $\v x\to\v y$ ($\v p\to\infty$ at fixed $\v q$) the leading terms $-4q^i\ln(\v p^2/\v q^2)/\v p^2$ of the two terms of $\mathcal R^i$ cancel and $\mathcal R^i=O(\ln p/p^3)$, so the kernel is bounded; at $\v x\to\v w$ ($\v q\to\infty$ at fixed $\v l=\v p-\v q$) the term $R_0^i(\v q,-\v l)/\v l^2\simeq2q^i\ln(\v l^2/\v q^2)/\v l^2$ gives $\ln|\v x-\v w|/|\v x-\v w|$. With $P\sim1/|\v p-\v w'|$: $\{p,\v w'\}$: $\lambda^{-1}$; $\{\v x,\v w,\v w'\}$: $\lambda^{-2}\ln\lambda$ ($d=4$); $\{\v x,\v y,\v w,\v w'\}$: $\lambda^{-2}$ ($d=6$). Integrable without the colour factor. For a local $n$ ($\v x=\v y$) the kernel vanishes, being antisymmetric under $\v x\leftrightarrow\v y$.

\noindent\textbf{Local two-point function, and local $o$} [\texttt{scan\_nlocal}]. If $n(p,r)$ contains $\delta^{(2)}(p-r)$, all three charges of a cubic structure coincide. All 329 clusters are absolutely integrable with the colour factor (margin $\ge0.90$). The kernel alone fails only for $\mathfrak K^{(z)}_\delta$ in all its contacts (with all charges at one point every such line contains $K(\v x,\v x,\v z)$), at $\{\v x,\v z\}$ and $\{\v x,\v y',\v z\}$, where the colour factors vanish as above. The same kernels at the triple coincidence enter the $(3\rho)$ equations for a local $o$ (Remark (iii)); there the colour factors $\Delta_{{\rm III}.R}(z)$ of Eq.~(Deltas) reduce at $\v x=\v y=\v x'$ to $-\mf G(x,z,w)$, $+\mf G(x,z,w)$ and $0$ for $R=$ III, V, VI, which vanish at $\v z\to\v x$ by the antisymmetry of $\mf G$; the $\Delta_{{\rm III}.R}(y')$ vanish at $\v x=\v y'=\v z$ in the same way. So the integrability at the triple coincidence asked for in Remark (iii) holds as well (by this power counting; the $\mf G$ factors themselves were not evaluated numerically here).

\noindent\textbf{Complex conjugate, angular averages.} The c.c.\ has the same modulus pointwise; the counting is the same. No angular average was used: all statements above are absolute.

\noindent\textbf{Large distance} [\texttt{ld\_pairs.py}]. The contact substitution acts on charge points only. The large-$\v z$ statements of Sec.~6.2 and Eq.~(eq:Mhatclaim) hold at arbitrary fixed charge positions, including coinciding ones, and the two members of a pair have identical ordered colour operators (Eq.~(eq:colid), operator identities; the pair members also have the same type of ordering: I.III and III.V both $\rho(x')\rho(x)\rho(y)$, I.III and III.VI after $x\leftrightarrow y$, I.V and III.III both $\rho\{\rho,\rho\}$), so their contact decompositions are identical and the pairs cancel line by line. In the Kz lines the colour factor does not depend on $\v w$, so (C3) applies unchanged. Numerically, with the non-uniform field (no limit of $U(\v z)$ at large $\v z$), for pair A ($\mathfrak K^{\rm III}_w$ lines with the $\mathfrak K_\Psi$ lines of III.V at $\v y'=\v w'$) and pair D (the $\mathcal E$ part of $\ell$ in the $\mathfrak K^{\rm III}_z$ lines with the $\mathfrak K_\Psi$ lines of III.VI, $x\leftrightarrow y$), in all three contacts: $\max\v z^2|\text{Group I member}|$ grows like $\ln|\v z|$ (e.g.\ $21.6,33.9,47.6,61.7$ for $|\v z|=10^{1\dots4}$, pair A, $y\to x$), while $\max\v z^2|\text{pair}|$ falls like $1/|\v z|$ ($5.1,0.61,0.12,0.012$). The merged charge point is confined by the support of $n$.

\subsection{(b) Do the contact terms cancel between rows? (Remark (i))}\label{sec:r3f-b}
\noindent [\texttt{classes.py}] Lines can cancel at the integrand level only if they have the same kernel after the contact substitution, the same integration variables and phase. The contact patterns $(x',x,y)\to(p,p,r)$, $(p,r,p)$, $(r,p,p)$ give different functions of $(p,r)$, also after $p\leftrightarrow r$ ($n$ symmetric), so lines with different contacts never share a kernel. Lines with the same kernel and contact exist (i) in Group~III, where Rows III.III, III.V, III.VI have the same kernels $\mathfrak K_{\widehat{\mathcal E}}$, $\mathfrak K_\Psi$, $\mathfrak K_\delta$, and (ii) in Rows II.I and III.II, which share $\mathfrak K_{\mathcal T}(w,w')$ (the c.c.\ of III.II is the $\mathfrak K_{\mathcal T}(w',w)$ class after $\v w\leftrightarrow\v w'$). The sums of the printed colour factors are
\begin{align}
\mathfrak K_{\mathcal T}(w,w')|_{y\to x}:&\ \ N_c\mf D(w,w')-\mf F(x,z,w')-N_c\mf D(x',w)+\mf F(x,x',z)\,,\notag\\
\mathfrak K_{\mathcal T}(w,w')|_{x'\to x}:&\ \ -N_c\mf D(w,w')+\mf F(x,z,w')-N_c\mf D(x,z)+N_c\mf D(x,w)\,,\notag\\
\mathfrak K_{y'}|_{y\to x}\ \text{(III.V+III.VI)}:&\ \ \pm\big[\mf F(x,x',w)-\mf F(x,y',w)-\mf F(x,x',z)+\mf F(x,y',z)\big]\,,
\end{align}
and eight-term combinations for $\mathfrak K_{y'}|_{x'\to x}$ and $\mathfrak K_{y'}|_{x'\to y}$ (III.III+III.V+III.VI); with random $SU(3)$ lines their moduli are $0.4$--$11$, comparable to the single lines ($1$--$15$). In II.I$+$III.II the quadrupoles $\mf Q$ cancel, but dipoles and $\mf F$ remain. So there is no integrand-level cancellation. A cancellation after integration would have to be an identity between different kernels for arbitrary $U$ and $n$; none is known and none is suggested by the structure. A cancellation against the contact terms of the three-$\rho$ evolution terms is impossible within the finite terms: those terms multiply $\ln(\vee/k^+)$ or $\ln(k^+/\Lambda)$ and are part of the BFKL and JIMWLK terms (they are the reordering commutators that complete the two-$\rho$ structures there), whereas the finite terms are the $\ln^0$ coefficient, Eq.~(eq:fin\_def). \textbf{Conclusion:} the contact terms are a genuine part of the result; their finiteness does not depend on a cancellation, Sec.~\ref{sec:r3f-a}.

\subsection{(c) The transition region of (C10)}\label{sec:r3f-c}
\noindent\textbf{Statement.} Rows II.I (second part), II.IV and III.III (second part) are, at fixed $t=p^+/k^+$, $|\mathcal A_B|^2$, $|\mathcal A_{CA}|^2$ and $-2{\rm Re}\,\mathcal A_{CA}^\dagger\mathcal A_B$, so their sum is $|\mathcal A_B-\mathcal A_{CA}|^2$, with the measured gluon Fourier transformed: $\int_{\v z}\widetilde{\mathcal A}(\v k,\v z)\widetilde{\mathcal A}^*(\v k,\v z)$, $\widetilde{\mathcal A}(\v k,\v z)=\int_{\v w}e^{-i\v k\cdot\v w}\mathcal A(\v w,\v z)$. Once $\v w$ (and $\v w'$) are integrated, there is no separate ``transition region'' in $\v z$: the question is the large-$|\v z|$ expansion of $\widetilde{\mathcal A}$, in which all configurations of $\v w$ (near the projectile, near $\v z$, and $|\v w|\sim|\v w-\v z|\sim|\v z|$) are included.

\noindent\textbf{Argument (method of regions).} For a target uniform outside a compact set (the setting of (C10)), $\mathcal A(\v w,\v z)$, $\v z=Z\hat{\v n}$, is a sum of products of homogeneous functions of point differences. In the outer region $\v w=Z\v u$ it is $Z^{-2}h(\v u)+O(Z^{-3})$ with $h$ smooth on $\mathbb R^2\setminus\{\text{singular points}\}$ and decaying at $|\v u|\to\infty$; the singular points are $\v u=0$ (measured gluon at the projectile: $\mathcal A_B$), $\v u=\hat{\v n}$ (pair: $\mathcal A_B$ and $\mathcal A_{CA}$) and $\v u=-t\hat{\v n}$ (parent at the projectile, collinear configuration: $\mathcal A_{CA}$). The Fourier transform at the large momentum $Z\v k$ of $h$ is determined by these singular points alone (integration by parts away from them gains $(kZ)^{-1}$ each time), and their contributions are exactly the Fourier transforms of the leading inner-region forms. Hence
\begin{align}
Z\,\widetilde{\mathcal A}(\v k,Z\hat{\v n})=c_{\rm soft}+c_{\rm joint}\,e^{-i\v k\cdot\v z}+c_{\rm coll}\,e^{+it\v k\cdot\v z}+O\Big(\frac{\ln Z}{Z}\Big)\,,
\end{align}
with $c_{\rm joint}$ the joint-limit coefficient of \texttt{fix\_tail} (computed at fixed $\v w-\v z$) and $c_{\rm soft}$, $c_{\rm coll}$ from the other two inner regions; the outer region contributes only at $O(Z^{-2})$. In $\mathcal A_B-\mathcal A_{CA}$, $c_{\rm joint}$ cancels (the parent pole, Eq.~(eq:dg\_bcA)). Squared: $|c_{\rm joint}|^2/Z^2$ is the joint tail (cancels in the sum of the three rows), $|c_{\rm soft}|^2/Z^2$ the soft tail of Row II.I (treated by its JIMWLK subtraction and the six-kernel sum $\mathfrak K_2$ at fixed $\v w,\v w'$), $|c_{\rm coll}|^2/Z^2$ the collinear region of Row II.IV (DGLAP), and the cross terms carry $e^{\pm i\v k\cdot\v z}$ or $e^{\pm i(1+t)\v k\cdot\v z}$, i.e.\ they are convergent Fourier integrals in $\v z$ (type (C4)). Equivalently, by Plancherel, $\int_{\v z}|\widetilde{\mathcal A}|^2=\int_{\v p}|\widetilde{\mathcal A}(\v k,\v p)|^2$, and the large-$\v z$ structure is the singular structure of $\widetilde{\mathcal A}(\v k,\v p)$ at $\v p=0$, $\v p=-\v k$ (soft parent) and $\v p=t\v k$ (collinear); no other momentum region corresponds to ``part of the pair far''.

\noindent\textbf{Numerical check} [\texttt{c10\_transition.py}, \texttt{c10\_extrap.py}]. Amplitudes of \texttt{fix\_tail} ($\mathcal A_B=N_BT_B/D_B$, $\mathcal A_{CA}=N_C\,C_{\rm amp}$, one charge at $\v x=0$, colour stripped, uniform far target), $\widetilde{\mathcal A}^{ij}(\v k,Z\hat{\v n})$ computed on the whole $\v w$ plane (Cartesian grid $h=0.1$--$0.15$ out to $|\v w|\approx10Z$, the $1/r$ singular parts at the three singular points subtracted with Gaussian cut-offs and added back analytically, a smooth window of radius $L$; the window error is $O(1/(kL))$, because the degree $-2$ tail does not oscillate for $\v w\perp\v k$, and was removed by extrapolation in $L=3Z,4.5Z,6Z$; $h$-independence checked). Residual $\max_{ij}|Z\widetilde{\mathcal A}-\text{prediction}|$, $\hat{\v n}$ at $0.6$ rad from $\v k$:
\begin{center}
\begin{tabular}{|c|c|cccc|}\hline
$t$ & amplitude ($|c_{\rm joint}|$) & $Z=10$ & $20$ & $40$ & $80$\\\hline
$1$ & $\mathcal A_B$ ($3.89$) & $7.5\times10^{-2}$ & $3.8\times10^{-2}$ & $2.8\times10^{-2}$ & $1.8\times10^{-2}$\\
$1$ & $\mathcal A_{CA}$ ($3.89$) & $5.9\times10^{-2}$ & $7.5\times10^{-2}$ & $4.2\times10^{-2}$ & $1.6\times10^{-2}$\\
$1$ & $\mathcal A_B-\mathcal A_{CA}$ ($0$) & $7.5\times10^{-2}$ & $3.8\times10^{-2}$ & $2.8\times10^{-2}$ & $1.8\times10^{-2}$\\
$0.3$ & $\mathcal A_B$ ($7.79$) & $0.41$ & $0.23$ & $0.12$ & $6.4\times10^{-2}$\\
$0.3$ & $\mathcal A_{CA}$ ($7.79$) & $0.81$ & $0.40$ & $0.15$ & $0.12$\\
$0.3$ & $\mathcal A_B-\mathcal A_{CA}$ ($0$) & $0.45$ & $0.22$ & $0.11$ & $5.5\times10^{-2}$\\\hline
\end{tabular}
\end{center}
A contribution of the outer (transition) region at the order of the tail would show up as a $Z$-independent residual of the size of $|c_{\rm joint}|$; the residuals fall like $\ln Z/Z$ ($t=1$: $0.075\to0.018$ from $Z=10$ to $80$, the ratio $0.24$ equals $(\ln80/80)/(\ln10/10)$) and like $1/Z$ ($t=0.3$: $Z\times$residual $=4.1,4.6,4.7$). The coherent amplitude has no $e^{-i\v k\cdot\v z}$ component at $O(1/Z)$. \textbf{The transition region is settled} for targets uniform at large distance. For a colour-neutral projectile and a non-uniform target, the outer region contains $U(\v w)$ at $|\v w|\sim Z$, which is not smooth on the scale $Z$; there only the paper's argument (the non-uniform remainder multiplies $\int_x\Kv(\v x-\v Y)\rho(\v x)=O(Z^{-2})$, so the outer region gives $O(Z^{-1})$ contributions to $\widetilde{\mathcal A}$ with random phases, whose square averages to $O(Z^{-4})$ for a target of finite correlation length) applies; it was not checked numerically.

\subsection{(d) Conditional convergence (Major 6)}\label{sec:r3f-d}
\noindent\textbf{Lemma (uniqueness of the measured-momentum Fourier transforms).} Let $f\in L^1_{\rm loc}(\mathbb R^2)$ be tempered (here: $|f(\v r)|\le C|\v r|^{-1}\ln|\v r|$ at large $|\v r|$), so that $\hat f$ is a tempered distribution, and let $\hat f$ be continuous in a neighbourhood of $\v q=\v k\neq0$. Then for every regulator $\phi_\epsilon(\v r)=\phi(\epsilon\v r)$ with $\phi(0)=1$ whose Fourier transform is integrable (Gaussian, $e^{-|\v r|}$, $e^{-c(\hat{\v r})|\v r|}$ with any positive continuous direction dependence $c$, or smooth cut-offs of any shape)
\begin{align}
\lim_{\epsilon\to0}\int d^2r\,e^{-i\v k\cdot\v r}\phi(\epsilon\v r)\,f(\v r)=\hat f(\v k)\,,
\end{align}
independently of $\phi$. Proof: the left side is $(\hat f*\hat\phi_\epsilon)(\v k)$ with $\hat\phi_\epsilon=\epsilon^{-2}\hat\phi(\cdot/\epsilon)$ an approximate identity; it converges to $\hat f(\v k)$ at points of continuity of $\hat f$ when $\hat f$ is locally integrable with at most polynomial growth. ``Radial integral first with $e^{-\epsilon r}$'', ``angular integral first'' and a Gaussian wave packet of the measured gluon are special cases. This is the precise sense in which the $\v w$, $\v w'$, $\v y'$, $\v y-\v y'$ (and, for the terms $\propto e^{-i\v k\cdot\v z}$, $\v z$) integrals at the measured momentum converge; at LO, $\int_we^{-i\v k\cdot\v w}\Kv(\v x-\v w)$, it is the same.

\noindent\textbf{Absolutely convergent grouping.} With the measured-momentum Fourier transforms understood in the sense of the Lemma (i.e.\ in momentum space for the measured gluon), all remaining integrals (positions of the unobserved gluon, of the charges, longitudinal fractions) are absolutely convergent in the grouping of the paper: pairs A--D under one $\v z$ integral; $\widehat{\mathcal M}$ with its $(\ell-\mathcal E)$ term; $\mathfrak K^{(y')}_\delta$ with $\mathfrak K^{(z)}_\delta$ (and the sub terms) under one $\v y'$ integral; the six $\mathfrak K_2$ together; the three rows of (C10) under one pair integral; the $(2\rho)$ and $(3\rho)$ lines in the same groups. The exceptions are the principal values $\mathcal P\frac1s$ (Sec.~\ref{sec:r3f-C7}) and the collinear point of Row I.V (Sec.~\ref{sec:r3f-IV}). The individual conditions:

\noindent\emph{(6.15), the two integrals of Eq.~(eq:sawid)} [\texttt{sawid\_reg.py}]. Both integrands are of the type of the Lemma (degree $-1$, the first with a logarithm; the second has a jump of its angular profile at $\alpha=\theta+\pi$, which makes $\hat f$ logarithmically singular only for $\cos\theta=0$, a set of measure zero in the subsequent $\v z$ integral). Each integral is therefore unique; ``angular integral first'' in the text refers to the subsequent $\v z$ integral of each term alone, $\int d^2z\,\chi(\theta)/\v z^2$, which is only conditionally convergent because $\chi$ has zero mean. In the grouping (the two terms at the same $\v z$) the two Fourier integrals cancel identically for every $\theta$, so nothing conditional is left. Numerically, with the isotropic Abel factor $e^{-\epsilon r}$ and two anisotropic ones, $e^{-\epsilon c(\alpha)r}$, $c=1+0.6\cos(2\alpha+0.3)$ and $c=1+0.5\cos(\alpha-1.1)$ (non-circular cut-off regions), at $\theta=0.7$: $|I_1-2\pi i\hat{\v k}\chi(\theta)/|\v k||=9.9\times10^{-4}$, $8.3\times10^{-4}$, $1.2\times10^{-3}$ and $|I_1+I_2|=6.5\times10^{-4}$, $1.05\times10^{-3}$, $8.1\times10^{-4}$ at $\epsilon=10^{-4}$, all falling $\propto\epsilon$ ($10\times$ larger at $\epsilon=10^{-3}$).

\noindent\emph{(C3)}. For fixed $\v z$, $\v w\mapsto\widehat{\mathcal M}^i(x,y,z,w)$ is locally integrable (Sec.~\ref{sec:r3f-a}: $\ln r/r$ at $\v w\to\v z,\v x$) with a degree $-1$ tail, so $\int_we^{-i\v k\cdot\v w}\widehat{\mathcal M}$ is unique by the Lemma; it is $O(\ln|z|/|z|^3)$ plus terms $\propto e^{-i\v k\cdot\v z}$ (Eq.~(eq:Mhatclaim)), and the subsequent $\v z$ integral is absolutely convergent (the oscillating terms are Fourier integrals in $\v z$, again unique). The opposite order does not exist: at fixed $\v w$ the angular mean of the $1/\v z^2$ tail is $\Kv^i(\v x-\v w)[\gamma_E+\ln(|\v k||\v x-\v w|)-\ln2]/\v z^2\neq0$ (the $\varphi$ part has zero mean, $c(\theta)$ has mean $-\ln2$), so $\int d^2z$ at fixed $\v w$ diverges logarithmically even after angular averaging. There is thus no second finite value that a different order could produce; (C3) is the definition of the measured momentum (the unobserved gluon is summed over at fixed $\v k$), not a choice. The tail cancellation needs a sharp $\v k$ only in the sense of the Lemma: any regulator of the measured gluon that is removed before the $\v z$ cut-off gives the same result.

\noindent\emph{(F.12), Eq.~(eq:app\_IVfin)} [\texttt{f12\_order.py}]. Here $\v y-\v y'$ is the Fourier variable (momentum $\v k/\xi$). Fourier first: the $\xi$ integrand is $O(\xi)$, absolutely convergent. $\xi$ first at fixed $\v r=\v y-\v y'$: $G(a)=\int_0^1d\xi\frac{1-\xi}{\xi}e^{-ia/\xi}=(1+ia)E_1(ia)-e^{-ia}$, $|G(a)|\le C(1+|\ln|a||)$ uniformly in a lower cut-off $\xi_0$, so by Fubini at finite $\xi_0$ (with a regulator of the $\v r$ integral) and dominated convergence for $\xi_0\to0$ the two orders agree. Toy with the structure of the term ($h=[\cos2\alpha+\ln\v r^2]e^{-\v r^2/2}$): order A $=-3.0246763985$, order B $=-3.0246763984$ (difference $5\times10^{-11}$). The ``conditional'' convergence at fixed coordinates is that of $\int du\,e^{-iau}/u$; it cannot shift the value.

\subsubsection{(C7): instantaneous terms}\label{sec:r3f-C7}
\noindent The defining order is: $\v y'$ integral (Fourier transform of the conjugate measured gluon at the intermediate momentum $(1+t)\v k$) at fixed $p^+$, then $p^+$. At fixed $t$ the bracket $h^m(\v r)=\Kv^m(\v r)[\Kv^j(\v x'-\v y')\mathcal L(\v y')-\Kv^j(\v x'-\v z)\mathcal L(\v z)]$, $\v r=\v y'-\v z$, is locally integrable with a degree $-1$ tail, so its Fourier transform $H((1+t)\v k)$ is unique (Lemma); because the bracket vanishes at $\v r\to0$, $H(\v q)=O(\v q^{-2})$ at large $\v q$ and $\int_0^\infty dt\,H((1+t)\v k)$ converges absolutely; this is the limit $\vee\to\infty$. The distributional form $\pi\delta(s)-i\mathcal P\frac1s$ is this quantity with the $t$ integral done first; then the perpendicular component $u$ of $\v r$ has no oscillating phase, and $\int du\,h$ converges only conditionally, through the odd, $s$-independent tail $-\Kv^\perp(\v r)\Kv^j(\v x'-\v z)\mathcal L(\v z)\simeq-c/u$. This is the only non-absolutely convergent piece, and it drops out for \emph{any} cut-off: a $\kappa$-independent change $\delta$ of the $u$ integral ($\kappa=\v k\cdot\v r$) contributes $\delta\int d\kappa\,e^{i\kappa}[\pi\delta(\kappa)+i\mathcal P\frac1\kappa]=\delta(\pi+i\cdot i\pi)=0$, and a cut-off varying with $\kappa$ on a scale $R\to\infty$ gives $O(1/R)$. Hence the symmetric prescription of (C7) is sufficient but not necessary, and no choice of the large-$u$ cut-off shifts the result. The coefficient functions of $\mathcal P\frac1s$ (after the $u$ integral) are H\"older continuous across $s=0$ except at isolated points where the merged charge lies on the line, where they have integrable logarithms, so the principal value exists (Sec.~\ref{sec:r3f-a}, $\mathcal P$ form). Toy check [\texttt{c7\_order.py}, $f$ Gaussian in place of $\Kv\mathcal L$]: C7TOY.

\subsubsection{The collinear point of Row I.V (not in the referee's list)}\label{sec:r3f-IV}
\noindent [\texttt{iv\_collinear.py}] At $\v z\to\v x$ the colour factor of $\mathfrak K^{\rm V}_{2a}$ vanishes ($N_c\mf D(x,x')-N_c\mf D(x,w')-\mf F(x,x',x)+\mf F(x,x,w')=0$), but that of $\mathfrak K^{\rm V}_{2b}$ does not. On circles $|\v z-\v x|=r$: $r^2\mathfrak K^{\rm V}_{2b}\to a(\varphi)$ with mean $3.5\times10^{-12}+5.4\times10^{-10}i$ at $r=10^{-4}$ (falling $\propto r^2$), first harmonic $2.2\times10^{-5}$ ($\propto r$) and second harmonic $0.0416$ (constant); without the $\frac{11}{6}$ term the mean is $0.0366$. So the point is a pure second harmonic $a_2\cos2(\varphi-\varphi_0)/r^2$: conditionally convergent, not absolutely. Its contribution is zero for any rotation-invariant cut-off around $\v z=\v x$ and finite but different for a non-circular one ($a_2\int d\varphi\cos2(\varphi-\varphi_0)\ln R(\varphi)$). The rotation-invariant choice is the one consistent with the $\frac{11}{6}$ subtraction, which comes from the dimensionally regularised UV/collinear integral. An absolutely convergent form is obtained by subtracting $\theta(R_0-|\v z-\v x|)\,a(\varphi)/|\v z-\v x|^2$, which integrates to zero with a rotation-invariant measure. Caveat: if the transverse indices are continued to $d_\perp=2-\epsilon$, the 2D-traceless tensor $\hat r^i\hat r^m-\frac12\delta^{im}$ has trace $O(\epsilon)$, which multiplies the $1/\epsilon$ of the collinear subtraction and gives a finite local term; whether it is present is the scheme question of Major~4 and is not decided here.

\subsection{Drop-in text for the paper}\label{sec:r3f-dropin}
{\begingroup\renewcommand{\eqref}[1]{(\texttt{\detokenize{#1}})}\renewcommand{\ref}[1]{\texttt{\detokenize{#1}}}\renewcommand{\label}[1]{}\providecommand{\pending}[1]{#1}
\noindent\textbf{1. Section~6.5 (sec:fin\_singlet), paragraph ``Remarks'' at its end: replace Remarks (i)--(iii)} (``\emph{(i) The contact terms of the cubic rows are part of the result \dots\ has not been examined.}'') by:\par\smallskip
\noindent\fbox{\parbox{0.97\textwidth}{\small
\textbf{Remarks.} (i) The contact terms of the cubic rows are part of the result for every colour-neutral projectile, whatever its three-point function $o$. They do not cancel among the rows at the level of the integrand: lines with the same kernel occur only in Rows~III.III, III.V and III.VI (kernels $\mathfrak K_{\widehat{\mathcal E}}$, $\mathfrak K_\Psi$, $\mathfrak K_\delta$ with the same contact) and in Rows~II.I and III.II (kernel $\mathfrak K_{\mathcal T}$), and the sums of their colour factors do not vanish (in II.I$+$III.II the quadrupoles $\mf Q$ cancel, dipoles and $\mf F$ remain). They cannot cancel against the contact terms of the three-$\rho$ evolution terms, which multiply $\ln(\vee/k^+)$ or $\ln(k^+/\Lambda)$. Their finiteness does not rely on a cancellation.
(ii) All contact kernels are integrable. We power-counted every line of Eqs.~\eqref{eq:sIIb}--\eqref{eq:sIIIVI} at every coincidence of its integration points (pairs and all larger clusters, in particular the coincidences of the merged charge point with $\v z$, $\v w$, $\v w'$, $\v y'$ and the other charge), with the colour factors as written: each line is absolutely integrable at each coincidence, with the singularities $\ln(1/r)/r$ or $1/r$ at pairs of points and a margin of at least $0.9$ powers of the scale at every larger cluster. Only one kernel is not integrable by itself, $\mathfrak K^{(z)}_\delta|_{x'\to x}$ of Rows~III.III, III.V and III.VI, which contains $K(\v x,\v x,\v z)=1/(\v x-\v z)^2$: at $\v z\to\v x$ (and at $\v x,\v y',\v z\to$ one point, where it combines with $\mathfrak K^{(y')}_\delta$) its colour factors, $-N_c\mf D(x,y)-N_c\mf D(z,w)+\mf F(x,y,z)+\mf F(x,z,w)$, $-N_c\mf D(x,y)+N_c\mf D(x,w)+\mf F(x,y,z)-\mf F(x,z,w)$ and $N_c\mf D(x,y)+N_c\mf D(y,z)-2\mf F(x,y,z)$, vanish identically by $\mf F(a,a,c)=N_c\mf D(a,c)$, so that the line is integrable for any target with H\"older continuous Wilson lines. In $\mathfrak K^{\rm III}_z|_{y\to x}$ and $\mathfrak K_{\mathcal T}|_{y\to x}$ the terms $1/(\v x-\v z)^2$ cancel within the kernel, leaving $\ln(1/r)/r$. The kernel of Row~I.II, $N_{2b}$, is homogeneous of degree $-1$ in $(\v x-\v w,\v y-\v w)$, bounded at $\v x\to\v y$ and $\propto\ln r/r$ at $\v x\to\v w$; with $P=\Kv(\v x-\v w')$ or $\Kv(\v y-\v w')$ it is integrable. The large-distance tails of the contact lines cancel pair by pair as at separate points, since the two members of each pair have the same ordered colour operator and hence the same contact decomposition. These statements were checked numerically with smooth, non-uniform $SU(3)$ Wilson-line fields at about $1200$ coincidence configurations.
(iii) The $(3\rho)$ equations assume $g_f=0$, as stated in (A2); a chiral part $f^{abc}g_f$ would add $t\,A\,\Phi\,g_f$ per cubic structure, also odd in $\v k$. For a local $o$ (or a local part of $n$) all three charge points coincide; the kernels remain integrable there, the only kernel that is not integrable by itself being again $\mathfrak K^{(z)}_\delta$ at $\v z\to\v x$, where the colour factors $\Delta_{{\rm III}.R}(z)$ reduce to $\mp\mf G(x,z,w)$ or $0$ and vanish by the antisymmetry of $\mf G$ (and the contact colour factors as in (ii)).}}\par\medskip

\noindent\textbf{2. Section~6.6 (sec:fin\_integrable), conditions.} Replace (C3) by:\par\smallskip
\noindent\fbox{\parbox{0.97\textwidth}{\small\begin{itemize}
\item[(C3)] In Eq.~\eqref{eq:GIIII} the $\v w$ integral, i.e.\ the Fourier transform at the measured momentum, is done before the $\v z$ integral. This is the definition of the measured momentum (the unobserved gluon is summed over at fixed $\v k$), not a choice: at fixed $\v w$ the brace has a $1/\v z^2$ tail whose angular mean, $\Kv^i(\v x-\v w)[\gamma_E+\ln(|\v k||\v x-\v w|)-\ln2]/\v z^2$, does not vanish, so the other order does not exist; after the $\v w$ integration, which is unique for any regulator of the measured gluon (Lemma below), the tail is absent in every direction, Eq.~\eqref{eq:sawid}, and the $\v z$ integral converges absolutely.\end{itemize}}}\par\smallskip
\noindent Replace (C7) by:\par\smallskip
\noindent\fbox{\parbox{0.97\textwidth}{\small\begin{itemize}
\item[(C7)] In the instantaneous terms of Group~III the $\v y'$ integral, the Fourier transform of the conjugate measured gluon at the intermediate momentum $(1+p^+/k^+)\v k$, is done before the $p^+$ integral (Section~\ref{sec:fin_inst}); this order defines the terms and gives an absolutely convergent $p^+$ integral. In the distributional form the component of $\v y'-\v z$ perpendicular to $\v k$ is integrated first; its only non-absolutely convergent part is independent of $s$ and is annihilated by $\int ds\,e^{-is}[\pi\delta(s)-i\mathcal P\frac1s]=0$, so that any cut-off at large distance, symmetric or not, gives the same result.\end{itemize}}}\par\smallskip
\noindent Replace (C10) by:\par\smallskip
\noindent\fbox{\parbox{0.97\textwidth}{\small\begin{itemize}
\item[(C10)] In the region in which both final-state gluons are far from the projectile together ($\v w,\v w',\v z\to\infty$ at fixed $\v w-\v z$), the finite terms of Row~II.IV, among them the remainder $N_{\mathcal R}$ of Eq.~\eqref{eq:dg_R} and the tensor term of Eq.~\eqref{eq:app_IVfin}, have the large-distance tail of Section~\ref{sec:dg_tail}. It cancels at the level of the amplitudes against the same region of the second part of Row~II.I and of Row~III.III, which therefore have to be integrated together with Row~II.IV under one integral over the position of the pair; the finite terms as written carry the compensating tail, $1\!:\!1\!:\!-2$ (Table~\ref{tab:tails}). The region in which only part of the pair is far does not contribute at this order: after the Fourier transform over the measured gluon, the large-$|\v z|$ expansion of each two-gluon amplitude is determined by the regions in which the measured gluon is near the projectile, near the other gluon, or in the collinear configuration, the remaining region being smooth (verified numerically on the full $\v w$ plane for $|\v z|$ up to $80/|\v k|$). For a coloured projectile and a target that is not uniform at large distance a remainder proportional to the total colour charge of the projectile is left; it is the non-singlet term of the JIMWLK formula and is absent for a colour-neutral projectile.\end{itemize}}}\par\smallskip
\noindent After the list of conditions, replace ``The checks are exact identities, pointwise tails in all directions and Fourier integrals at fixed values of the other coordinates; a complete numerical integration \dots\ has not been carried out.'' by:\par\smallskip
\noindent\fbox{\parbox{0.97\textwidth}{\small
\textbf{Convergence.} All integrals over a coordinate that carries the measured phase ($\v w$, $\v w'$, $\v y'$, $\v y-\v y'$, and $\v z$ in the terms $\propto e^{-i\v k\cdot\v z}$) are Fourier integrals of locally integrable, tempered functions whose transforms are continuous at the measured momentum. For such a function $f$ and any regulator $\phi(\epsilon\v r)$ with $\phi(0)=1$ and integrable Fourier transform (an exponential or Gaussian factor, of any direction dependence, or a smooth cut-off of any shape), $\lim_{\epsilon\to0}\int d^2r\,e^{-i\v k\cdot\v r}\phi(\epsilon\v r)f(\v r)=\hat f(\v k)$, since the left-hand side is the convolution of $\hat f$ with an approximate identity. These integrals, like $\int_we^{-i\v k\cdot\v w}\Kv(\v x-\v w)$ at LO, are therefore unique and independent of the order of their radial and angular parts. With them done first, every remaining integral, over the positions of the unobserved gluon and of the charges and over the longitudinal fractions, converges absolutely in the grouping of this section (pairs A--D; $\widehat{\mathcal M}$; $\mathfrak K^{(y')}_\delta$ with $\mathfrak K^{(z)}_\delta$; the six $\mathfrak K_2$; the three rows of (C10)), except for the principal values $\mathcal P\frac1s$ and the collinear point $\v z\to\v x$ of Row~I.V, which is conditionally convergent (a pure second harmonic in the angle) and is defined with a rotation-invariant cut-off, as the $\frac{11}{6}$ subtraction that accompanies it. The two integrals of Eq.~\eqref{eq:sawid} and the two orders of the $\xi$ and $\v y-\v y'$ integrals in Eq.~\eqref{eq:app_IVfin} were evaluated with different regulators and orders, with the same result. The checks are exact identities, power counting at all coincidences, pointwise tails in all directions and Fourier integrals at fixed values of the other coordinates; a complete numerical integration of the finite terms over all transverse positions with Wilson lines has not been carried out.}}\par\medskip

\noindent\textbf{3. Eq.~(eq:sawid), last sentence of item 3} (``Since $\chi$ has zero angular mean, each of them alone is only conditionally convergent (angular integral first), and not convergent at all if $U(\v z)$ depends on the direction of $\v z$.'') $\to$\par\smallskip
\noindent\fbox{\parbox{0.97\textwidth}{\small Each of the two Fourier integrals is unique (independent of the regulator and of the order of the radial and angular integrals, see the paragraph ``Convergence'' below; checked with isotropic and anisotropic exponential regulators); what is only conditionally convergent is the subsequent $\v z$ integral of each term alone, since $\chi$ has zero angular mean, and it would not converge if $U(\v z)$ depended on the direction of $\v z$. The two terms cancel for every $\theta$, and in their sum nothing conditional is left.}}\par\medskip

\noindent\textbf{4. Appendix F, after Eq.~(eq:app\_IVfin)}: replace ``at fixed coordinates the $\xi$ integral converges conditionally, with an integrable logarithm on the line $\v k\cdot(\v y-\v y')=0$'' by ``at fixed coordinates the $\xi$ integral converges conditionally, $\int_0^1d\xi\frac{1-\xi}{\xi}e^{-ia/\xi}=(1+ia)E_1(ia)-e^{-ia}$, with an integrable logarithm on the line $\v k\cdot(\v y-\v y')=0$, and gives the same result as the order with the $\v y-\v y'$ integral first''. In Appendix F, Row~I.V, after ``\dots is integrable with the angular integral done first, by virtue of the $\frac{11}{6}$ term'' add: ``(the remaining singularity is a pure second harmonic in the angle of $\v z-\v x$, i.e.\ conditionally convergent; the angular integral first, a rotation-invariant cut-off, is the regularisation in which the $\frac{11}{6}$ subtraction was obtained)''.\par\medskip

\noindent\textbf{5. Section~6.8 (summary), last sentence}, ``the term-by-term check of this cancellation on the finite terms as written is still open, (C10)'' $\to$ ``the cancellation holds term by term on the finite terms as written, including the region in which only part of the pair is far, (C10)''. \textbf{Section~5.4}, last sentence ``Not verified are the transition region in which only part of the pair is far, \dots'' $\to$ ``The region in which only part of the pair is far does not contribute at this order (Section~\ref{sec:fin_integrable}, (C10)). Not verified are the interchange of the $p^+$ integral with the joint limit in Row~III.III (checked numerically for Row~II.I only) and a full transverse integration with Wilson lines.''\par\medskip

\noindent\textbf{6. Conclusions}, open points (i), (ii) $\to$ ``(i) The large-distance tail of the collinear row at large position of the parent gluon cancels against the corresponding region of two other rows, provided the three rows are integrated together over the position of the gluon pair; the region in which only part of the pair is far does not contribute. (ii) All contact kernels are integrable at all coincidences; the only kernel that is not integrable by itself, $K(\v x,\v x,\v z)=1/(\v x-\v z)^2$ in the instantaneous terms, comes with a colour factor that vanishes at $\v z=\v x$. The contact terms do not cancel among the rows.''
\par\endgroup}

\subsection{What is not done / open}\label{sec:r3f-open}
\noindent (1) The power counting is numerical (exponent fits, $\lambda\ge10^{-5}$, two smooth fields, three base configurations per line) plus the analytic arguments given; it is not a proof for every configuration (e.g.\ collinear or otherwise degenerate base configurations were not sampled specifically; the $\arccos/\sqrt{\cdot}$ structure of $\mathcal A_0$ on the collinear line $x,z,w$ was taken from the paper). (2) The quadratic rows (I.V second part, II.I second part, II.IV, III.III second part) were not power-counted here, except the collinear point of Row~I.V; their integrability is as stated in the paper. (3) Row I.II $N_{2b}$: analytic estimate only (no coordinate form). (4) Pairs B and C of the soft block were not re-checked numerically in contact form (operator argument only). (5) (C10): the numerical check is for one charge, colour stripped (uniform far target), two values of $t$, one direction; the non-uniform target with a colour-neutral projectile rests on the paper's argument. Uniformity of the transition estimate in $t$ near $t\to0,\infty$ (where the BFKL and JIMWLK subtractions act) was not examined beyond \texttt{fix\_tail}. (6) The scheme dependence of the collinear point of Row I.V ($d_\perp$ continuation of the tensor) belongs to Major~4 and is open. (7) The $(3\rho)$ colour factors $\mf G$ at the triple coincidence were reduced analytically, not evaluated numerically.
